\documentclass[11pt,a4paper]{article}

\usepackage[margin=2.2cm]{geometry}
\usepackage[indonesian]{babel}
\usepackage{lmodern}
\usepackage[T1]{fontenc}
\usepackage{xcolor}
\usepackage{hyperref}
\usepackage{enumitem}
\usepackage{tcolorbox}
\usepackage{listings}
\usepackage{titlesec}
\usepackage{parskip}
\usepackage{fancyhdr}

\definecolor{Navy}{HTML}{141F36}
\definecolor{Blue}{HTML}{3166DC}
\definecolor{Mint}{HTML}{43BEA4}
\definecolor{Light}{HTML}{F5F7FB}
\definecolor{CodeBG}{HTML}{1C212B}
\definecolor{Gray}{HTML}{697180}

\hypersetup{
    colorlinks=true,
    linkcolor=Blue,
    urlcolor=Blue
}

\pagestyle{fancy}
\fancyhf{}
\lhead{\textcolor{Navy}{DAWET}}
\rhead{\textcolor{Gray}{Meeting 02}}
\cfoot{\thepage}

\titleformat{\section}{\Large\bfseries\color{Navy}}{}{0em}{}
\titleformat{\subsection}{\large\bfseries\color{Blue}}{}{0em}{}

\lstdefinestyle{htmlstyle}{
    language=HTML,
    basicstyle=\ttfamily\small\color{white},
    backgroundcolor=\color{CodeBG},
    keywordstyle=\color{Mint},
    commentstyle=\color{Gray},
    stringstyle=\color{yellow},
    breaklines=true,
    frame=none,
    showstringspaces=false,
    columns=fullflexible
}

\newtcolorbox{conceptbox}{
    colback=Light,
    colframe=Blue,
    boxrule=0.7pt,
    arc=3mm,
    left=4mm,right=4mm,top=2mm,bottom=2mm
}

\begin{document}

\begin{titlepage}
    \centering
    \vspace*{3cm}
    {\Huge\bfseries\color{Navy} DAWET\par}
    \vspace{0.5cm}
    {\Huge\bfseries\color{Blue} HTML dari Nol\par}
    \vspace{0.7cm}
    {\Large Meeting 02 — Modul Belajar Pemula\par}
    \vspace{1.5cm}
    \begin{tcolorbox}[colback=Light,colframe=Mint,arc=3mm]
        \centering
        Dari struktur HTML dasar menuju fondasi Portal UMKM Kampus.
    \end{tcolorbox}
    \vfill
    {\large Study Club Web Development\par}
\end{titlepage}

\tableofcontents
\newpage

\section{Tujuan Pembelajaran}

Setelah menyelesaikan materi ini, peserta diharapkan mampu:

\begin{enumerate}[label=\arabic*.]
    \item Menjelaskan fungsi HTML dalam website.
    \item Membedakan tag, element, dan attribute.
    \item Membuat struktur dasar dokumen HTML.
    \item Menggunakan heading, paragraph, list, link, dan image.
    \item Membuat halaman profil UMKM sederhana.
\end{enumerate}

\section{Apa Itu HTML?}

HTML adalah singkatan dari \textit{HyperText Markup Language}. HTML digunakan untuk menyusun struktur dan memberi makna pada konten halaman web.

HTML bukan bahasa pemrograman seperti JavaScript atau Python. HTML merupakan \textit{markup language}, yaitu bahasa penanda yang memberi tahu browser tentang struktur sebuah dokumen.

\begin{conceptbox}
\textbf{Analogi sederhana}

\begin{itemize}
    \item HTML = struktur dan ruangan sebuah rumah.
    \item CSS = warna, dekorasi, dan tampilan rumah.
    \item JavaScript = interaksi dan perilaku di dalam rumah.
\end{itemize}
\end{conceptbox}

\section{Tag, Element, dan Attribute}

Perhatikan contoh berikut:

\begin{lstlisting}[style=htmlstyle]
<a href="https://example.com">Kunjungi website</a>
\end{lstlisting}

\subsection{Tag}

Tag adalah penanda yang ditulis menggunakan tanda kurung sudut. Contohnya \texttt{<a>} sebagai tag pembuka dan \texttt{</a>} sebagai tag penutup.

\subsection{Element}

Element adalah keseluruhan bagian HTML yang terdiri dari tag pembuka, isi, dan tag penutup.

\subsection{Attribute}

Attribute memberikan informasi tambahan pada sebuah tag. Pada contoh sebelumnya, \texttt{href} adalah attribute yang berisi alamat tujuan link.

\section{Struktur Dasar HTML}

Gunakan struktur berikut sebagai titik awal dokumen HTML:

\begin{lstlisting}[style=htmlstyle]
<!DOCTYPE html>
<html lang="id">
<head>
    <meta charset="UTF-8">
    <title>Halaman Saya</title>
</head>
<body>
    <h1>Hello, DAWET!</h1>
    <p>Ini halaman web pertama saya.</p>
</body>
</html>
\end{lstlisting}

\subsection{\texttt{<!DOCTYPE html>}}

Deklarasi ini memberi tahu browser bahwa dokumen menggunakan HTML5.

\subsection{\texttt{<html>}}

Element utama yang membungkus seluruh dokumen HTML. Attribute \texttt{lang="id"} menunjukkan bahwa bahasa dokumen adalah bahasa Indonesia.

\subsection{\texttt{<head>}}

Bagian yang berisi informasi halaman, seperti judul, metadata, dan referensi file CSS.

\subsection{\texttt{<body>}}

Bagian yang berisi konten utama yang ditampilkan di halaman browser.

\section{Heading dan Paragraph}

Heading digunakan untuk judul atau subjudul. HTML menyediakan heading dari \texttt{h1} sampai \texttt{h6}.

\begin{lstlisting}[style=htmlstyle]
<h1>Warung Berkah</h1>
<h2>Menu Favorit</h2>
<p>Kami menyediakan makanan rumahan
dengan harga bersahabat.</p>
\end{lstlisting}

Gunakan heading berdasarkan struktur dan tingkat kepentingan konten, bukan hanya untuk memperbesar ukuran teks.

\section{List}

HTML menyediakan dua jenis list yang umum digunakan.

\subsection{Unordered List}

\texttt{<ul>} digunakan untuk daftar yang tidak memiliki urutan khusus.

\begin{lstlisting}[style=htmlstyle]
<ul>
    <li>Nasi Goreng</li>
    <li>Mie Ayam</li>
    <li>Es Teh</li>
</ul>
\end{lstlisting}

\subsection{Ordered List}

\texttt{<ol>} digunakan untuk daftar yang memiliki urutan.

\begin{lstlisting}[style=htmlstyle]
<ol>
    <li>Pilih menu</li>
    <li>Hubungi penjual</li>
    <li>Konfirmasi pesanan</li>
</ol>
\end{lstlisting}

Element \texttt{<li>} digunakan untuk setiap item dalam kedua jenis list tersebut.

\section{Link dengan \texttt{<a>}}

Element \texttt{<a>} digunakan untuk membuat link. Attribute \texttt{href} menentukan tujuan link.

\begin{lstlisting}[style=htmlstyle]
<a href="https://instagram.com/">
    Instagram UMKM
</a>
\end{lstlisting}

Link dapat digunakan untuk:

\begin{itemize}
    \item Mengarahkan pengunjung ke Instagram.
    \item Mengarahkan pengunjung ke WhatsApp.
    \item Membuka halaman lain.
    \item Mengarah ke bagian tertentu dalam halaman.
\end{itemize}

\section{Gambar dengan \texttt{<img>}}

Element \texttt{<img>} digunakan untuk menampilkan gambar.

\begin{lstlisting}[style=htmlstyle]
<img src="images/menu.jpg"
     alt="Foto menu makanan">
\end{lstlisting}

\begin{itemize}
    \item \texttt{src}: lokasi atau alamat file gambar.
    \item \texttt{alt}: teks alternatif yang menjelaskan gambar.
\end{itemize}

Attribute \texttt{alt} penting untuk aksesibilitas dan membantu ketika gambar gagal dimuat.

\section{Contoh Proyek: Profil UMKM}

Buat file bernama \texttt{index.html}, lalu isi dengan contoh berikut:

\begin{lstlisting}[style=htmlstyle]
<!DOCTYPE html>
<html lang="id">
<head>
    <meta charset="UTF-8">
    <title>Warung Berkah</title>
</head>
<body>
    <h1>Warung Berkah</h1>
    <p>Makanan rumahan untuk mahasiswa.</p>

    <h2>Menu</h2>
    <ul>
        <li>Nasi Goreng — Rp12.000</li>
        <li>Es Teh — Rp4.000</li>
    </ul>

    <a href="https://wa.me/628123456789">
        Hubungi via WhatsApp
    </a>
</body>
</html>
\end{lstlisting}

Cobalah mengganti nama UMKM, deskripsi, menu, harga, dan link kontak dengan data UMKM pilihanmu.

\section{Praktik Mandiri}

\subsection{Persiapan Folder}

Buat struktur folder berikut:

\begin{lstlisting}[style=htmlstyle]
DAWET/
└── minggu-02/
    └── index.html
\end{lstlisting}

\subsection{Tugas Utama}

Buat halaman profil UMKM dengan ketentuan:

\begin{enumerate}[label=\arabic*.]
    \item Memiliki satu judul utama menggunakan \texttt{h1}.
    \item Memiliki deskripsi singkat menggunakan \texttt{p}.
    \item Memiliki minimal tiga menu menggunakan \texttt{ul} dan \texttt{li}.
    \item Memiliki satu gambar menggunakan \texttt{img} dan \texttt{alt}.
    \item Memiliki satu link kontak menggunakan \texttt{a} dan \texttt{href}.
    \item Memiliki struktur \texttt{html}, \texttt{head}, dan \texttt{body}.
\end{enumerate}

\subsection{Tantangan Tambahan}

Jika tugas utama sudah selesai, tambahkan:

\begin{itemize}
    \item Link Instagram.
    \item Daftar langkah pemesanan menggunakan \texttt{ol}.
    \item Harga untuk setiap menu.
    \item Komentar HTML untuk menandai bagian kode.
    \item Judul halaman yang sesuai pada element \texttt{title}.
\end{itemize}

\section{Checkpoint Pemahaman}

Coba jawab pertanyaan berikut menggunakan bahasamu sendiri:

\begin{enumerate}[label=\arabic*.]
    \item Apa fungsi HTML?
    \item Apa perbedaan tag, element, dan attribute?
    \item Apa fungsi \texttt{href} dan \texttt{src}?
    \item Kapan kita menggunakan \texttt{ul} dan kapan menggunakan \texttt{ol}?
    \item Apa perbedaan \texttt{head} dan \texttt{body}?
    \item Mengapa gambar sebaiknya memiliki attribute \texttt{alt}?
\end{enumerate}

\section{Tugas GitHub dan README}

Simpan proyek pada repository GitHub pribadi. Tambahkan README sederhana yang berisi:

\begin{itemize}
    \item Nama dan identitas.
    \item Judul proyek.
    \item Deskripsi singkat.
    \item Teknologi yang digunakan.
    \item Hal yang dipelajari.
    \item Kendala dan solusi.
    \item Screenshot hasil halaman.
\end{itemize}

\section{Referensi}

\begin{enumerate}
    \item MDN Web Docs — Basic HTML Syntax:
    \url{https://developer.mozilla.org/en-US/docs/Learn_web_development/Core/Structuring_content/Basic_HTML_syntax}
    \item MDN Web Docs — HTML Elements Reference:
    \url{https://developer.mozilla.org/en-US/docs/Web/HTML/Reference/Elements}
    \item W3C Markup Validation Service:
    \url{https://validator.w3.org/}
\end{enumerate}

\section{Penutup}

Jangan berfokus pada menghafal semua tag HTML sekaligus. Fokus utama pada sesi ini adalah memahami struktur dokumen dan mampu membuat halaman sederhana dari nol.

\begin{conceptbox}
\centering
\textbf{Satu halaman hari ini. Satu website besok.}

Mulai dari struktur kecil, bangun proyek yang nyata.
\end{conceptbox}

\end{document}
